Per resoldre el segon apartat, calia saber la tensió eficaç o màxima del generador. Com que aquesta dada no apareix a l'enunciat, la qualificació de l'apartat (a) és de 2,5.

\begin{solapartat}
\punts{0,5} A partir del gràfic podem veure que el període és: $T = 0,02\ \mathrm{s}$

\punts{0,6} Llavors, la freqüència és:
\[ f = \frac{1}{T} = \frac{1}{0,02} = 50,0\ \mathrm{Hz} \]

\punts{0,6} A partir del gràfic podem veure que la intensitat màxima és:
\[ I_\textit{màx} = 3,0\ \mathrm{A} \]

\punts{0,2} La freqüència angular és:
\[ \omega = 2\pi f = 100\pi\ \mathrm{rad/s} \]

\destaca{Primera opció}

\punts{0,2} $I(t) = I_\textit{màx}\cos(\omega t + \varphi_0)$

\punts{0,2} Inicialment, $I(t = 0) = 0$
\[ 0 = \cos(\varphi_0) \;\Rightarrow\; \varphi_0 = \mathrm{ArcCos}(0) = \pi/2\ \mathrm{rad}. \]

\punts{0,2} Finalment, l'equació és:
\[ I(t) = 3\cos\left(100\pi\,t + \frac{\pi}{2}\right),\ I \text{ en A i } t \text{ en s}. \]

\destaca{Alternativament:}

\punts{0,2} $I(t) = I_\textit{màx}\sin(\omega t + \varphi_0)$

\punts{0,2} Inicialment, $I(t = 0) = 0$
\[ 0 = \sin(\varphi_0) \;\Rightarrow\; \varphi_0 = \mathrm{ArcSin}(0) = 0. \]

\punts{0,2} Finalment, l'equació és:
\[ I(t) = 3\sin(100\pi\,t),\ I \text{ en A i } t \text{ en s}. \]
\end{solapartat}

\begin{solapartat}
% DUBTE: la pauta oficial (pau_fisi22sp.pdf, pàg. 11) indica explícitament que aquest
% apartat no es pot resoldre perquè l'enunciat no dona la tensió eficaç ni la màxima
% del generador (només la intensitat), i que per aquest motiu tota la qualificació del
% problema (2,5 punts) es concentra en l'apartat (a). No hi ha, doncs, cap solució
% numèrica ni criteri de correcció per a aquest apartat b) a la pauta.
\textit{La pauta oficial no dona cap solució per a aquest apartat: indica que l'enunciat no aporta la dada (tensió eficaç o màxima del generador) necessària per calcular-lo, i que per això tots els 2,5 punts del problema es qualifiquen amb l'apartat (a).}
\end{solapartat}
